\section*{Capítulo \ref{ch:b3:ocean-dynamics} --- Dinâmica dos oceanos}

\begin{solution}{exo:b3:ocean-dynamics:1}
$f = \qty{9.37e-5}{s^{-1}}$; $\tau/(\rho f) = 0.12/(1025\times\num{9.37e-5})
= \qty{1.25}{m^2/s}$, à direita do vento: sul. Através de
\qty{1000}{km}, \qty{1.25e6}{m^3/s}, \qty{1.25}{\sverdrup}.
\end{solution}

\begin{solution}{exo:b3:ocean-dynamics:2}
$f = \qty{1.031e-4}{s^{-1}}$; $D_E = \pi\sqrt{2\times0.05/\num{1.031e-4}}
= \qty{98}{m}$. A corrente ali é $\mathrm e^{-\pi} = \num{0.043}$ da
corrente superficial, fluindo no sentido oposto.
\end{solution}

\begin{solution}{exo:b3:ocean-dynamics:3}
$\beta = 2\Omega\cos\ang{30}/a = \qty{1.98e-11}{m^{-1}.s^{-1}}$; $V =
\num{-1e-7}/(1025\times\num{1.98e-11}) = \qty{-4.92}{m^2/s}$; através de
\qty{5000}{km}, \qty{-24.6}{\sverdrup}: para o sul.
\end{solution}

\begin{solution}{exo:b3:ocean-dynamics:4}
As duas \omterm{def:b3:ocean-dynamics:tides}{constituintes de maré} voltam a ficar em fase após $1/(1/12 - 1/12.4206) =
\qty{354}{h}$, \num{14.77}\text{ dias}, metade de um mês sinódico.
\end{solution}

\begin{solution}{exo:b3:ocean-dynamics:5}
Em \ang{40} N, $f = \qty{9.37e-5}{s^{-1}}$, $\beta =
\qty{1.75e-11}{m^{-1}.s^{-1}}$, $\tau_x = \qty{0.05}{Pa}$ e
$\partial\tau_x/\partial y = 0.1/\num{1.11e6} = \qty{9.0e-8}{Pa/m}$. O
primeiro termo, $-\num{9.0e-8}/(1025\times\num{9.37e-5}) =
\qty{-9.38e-7}{m/s}$; o segundo, $0.05\times\num{1.75e-11}/(1025\times
(\num{9.37e-5})^2) = \qty{+9.7e-8}{m/s}$. $w_E = \qty{-8.4e-7}{m/s}$,
\qty{-26.5}{m/yr}: a água é empurrada para baixo.
\end{solution}

\begin{solution}{exo:b3:ocean-dynamics:6}
$\delta_S = \num{2e-6}/\num{2e-11} = \qty{100}{km}$; velocidade $=
\num{30e6}/(\num{1e5}\times800) = \qty{0.375}{m/s}$.
\end{solution}

\begin{solution}{exo:b3:ocean-dynamics:7}
$\beta = \qty{1.98e-11}{m^{-1}.s^{-1}}$; $\delta_M = (10^4/\num{1.98e-11})^{1/3}
= \qty{79.6}{km}$, e o jato tem $3.628\,\delta_M = \qty{289}{km}$ de largura.
Com $A_H$ oito vezes menor, ambas são reduzidas à metade, já que $8^{1/3} = 2$.
\end{solution}

\begin{solution}{exo:b3:ocean-dynamics:8}
$f = \qty{8.57e-5}{s^{-1}}$; $v = (21.3 - 14.5)/(\num{8.57e-5}\times\num{8e4})
= \qty{0.99}{m/s}$, positivo com $B$ à direita: olhando de $A$ para $B$,
a leste, a corrente segue para o norte. A superfície do mar fica
$6.8/9.81 = \qty{0.69}{m}$ mais alta em $B$; a coluna de água em $B$ é a
mais leve e mais quente, à direita da corrente.
\end{solution}

\begin{solution}{exo:b3:ocean-dynamics:9}
Amplitude de sizígia $1.5K(1 + 0.46) = \qty{0.78}{m}$; amplitude de quadratura $1.5K(1 - 0.46)
= \qty{0.29}{m}$; razão $1.46/0.54 = 2.7$.
\end{solution}

\begin{solution}{exo:b3:ocean-dynamics:10}
Como no \cref{thm:b3:ocean-dynamics:kelvin}, $u = \pm(g/c)\eta$ para $F(x \mp
ct)$, e $fu = -g\,\partial\eta/\partial y$ torna-se $\beta_0y(\pm g/c)G =
-gG'$, logo $G = \exp(\mp\beta_0y^2/2c)$. Apenas o sinal superior, correspondente à
onda para leste, é limitado nos dois lados do equador. $\beta_0 =
\qty{2.29e-11}{m^{-1}.s^{-1}}$, $\sqrt{2.5/\num{2.29e-11}} = \qty{330}{km}$.
\end{solution}

\begin{solution}{exo:b3:ocean-dynamics:11}
(a)~Integrando $\beta v = f\,\partial w/\partial z$ desde o fundo, $w = 0$,
até o topo da camada abissal, $w = W_0$: $\beta V = fW_0$, e $f/\beta =
a\tan\varphi$. (b)~$V = \num{6.371e6}\times0.577\times10^{-7} =
\qty{0.368}{m^2/s}$, \qty{1.84}{\sverdrup} através de \qty{5000}{km}, em direção ao polo.
(c)~A partir de uma corrente de contorno oeste profunda, alimentada pelas fontes de altas latitudes;
o interior escoa em direção às fontes, enquanto a corrente de contorno se afasta
delas.
\end{solution}

\begin{solution}{exo:b3:ocean-dynamics:12}
(a)~Com $z$ para cima a partir do fundo, $A_zW'' = \mathrm if(W - W_g)$, o
gradiente de pressão é o do escoamento geostrófico acima. $W - W_g =
C\mathrm e^{-(1+\mathrm i)z/\delta}$ é limitada para cima, e $W(0) = 0$ dá
$C = -W_g$. (b)~O déficit de transporte é $\int_0^\infty(W - W_g)\,\mathrm dz
= -W_g\delta/(1+\mathrm i) = -\tfrac12\delta(1 - \mathrm i)W_g$, ou seja
$-\tfrac12\delta(u_g + v_g,\ v_g - u_g)$. Como o escoamento geostrófico é
não divergente, $w = -\nabla\cdot(\text{déficit}) = \tfrac12\delta(\partial_x
v_g - \partial_y u_g) = \tfrac12\delta\zeta_g$. (c)~$\delta =
\sqrt{0.02/10^{-4}} = \qty{14.1}{m}$; $r = 10^{-4}\times14.1/8000 =
\qty{1.77e-7}{s^{-1}}$, um tempo de \textit{spin-down} de \num{65}\text{ dias}.
\end{solution}

\begin{solution}{pb:b3:ocean-dynamics:1}
\textbf{1.} $\tau_x = \qty{-0.1}{Pa}$ (de leste), 0 e \qty{+0.1}{Pa}
(de oeste).

\textbf{2.} $f = \num{4.99e-5}$, \num{8.37e-5} e
\qty{1.117e-4}{s^{-1}}.

\textbf{3.} $V_E = -\tau_x/(\rho f)$: \qty{+1.96}{m^2/s} (norte) em
\ang{20} N e \qty{-0.87}{m^2/s} (sul) em \ang{50} N. Elas convergem.

\textbf{4.} $D_E = \pi\sqrt{2\times0.05/\num{8.37e-5}} = \qty{109}{m}$.

\textbf{5.} $\tau_x = 0$ em \ang{35} N, portanto $w_E = -(1/\rho f)\,
\partial\tau_x/\partial y = -\pi\tau_0/(\rho fL) = \qty{-1.10e-6}{m/s}$,
\qty{-34.7}{m/yr}, para baixo.

\textbf{6.} $500/34.7 = \num{14.4}\text{ anos}$.

\textbf{7.} $L = a\times\pi/6 = \qty{3336}{km}$; $\beta =
\qty{1.875e-11}{m^{-1}.s^{-1}}$.

\textbf{8.} $\operatorname{curl}_z\vec\tau = -\pi\tau_0/L =
\qty{-9.42e-8}{Pa/m}$.

\textbf{9.} $V = \num{-9.42e-8}/(1025\times\num{1.875e-11}) =
\qty{-4.90}{m^2/s}$, para o sul.

\textbf{10.} $4.90\times\num{6e6} = \qty{2.94e7}{m^3/s}$,
\qty{29.4}{\sverdrup} para o sul.

\textbf{11.} $\qty{4.9}{mm/s}$; \qty{1000}{km} em \qty{2.0e8}{s},
\num{6.5}~anos.

\textbf{12.} $fw_E/\beta = \num{8.37e-5}\times(\num{-1.10e-6})/
\num{1.875e-11} = \qty{-4.90}{m^2/s}$: o mesmo, porque o
transporte de Ekman se anula onde $\tau_x = 0$. Uma coluna comprimida à taxa $w_E$
mantém sua \omterm{def:b3:atmospheric-dynamics:pv}{vorticidade potencial} $f/h$ apenas deslocando-se para valores menores de $f$, para o sul,
à taxa $\beta V_g = fw_E$.

\textbf{13.} Uma camada limite só pode ser aquela em que a correção para
o interior decai afastando-se da parede; com $\beta > 0$ isso é
$\mathrm e^{-x/\delta_S}$, que decai para leste, a partir da parede oeste
apenas. Em termos de vorticidade, o atrito em uma corrente para o norte contra uma
parede leste removeria vorticidade positiva, o sinal errado.

\textbf{14.} $\delta_S = \num{e-6}/\num{1.875e-11} = \qty{53.3}{km}$.

\textbf{15.} $A = \qty{4.90}{m^2/s}$, $X/\delta_S = 112.5$:
$4.90\times111.5 = \qty{546}{m^2/s}$, \qty{0.55}{m/s} através de
\qty{1000}{m}.

\textbf{16.} Onde $(X/\delta_S)\mathrm e^{-x/\delta_S} = 1$: $x =
\delta_S\ln(X/\delta_S) = \qty{252}{km}$.

\textbf{17.} $A(X - \delta_S - x_0) = 4.90\times(6000 - 53 - 252)\times
10^3 = \qty{27.9}{\sverdrup}$. O escoamento interior ainda segue para o sul na
faixa junto à parede; o déficit, $A(x_0 + \delta_S) =
\qty{1.5}{\sverdrup}$, é o transporte do interior que teria atravessado
essa faixa.

\textbf{18.} $\delta_M = (\num{5e3}/\num{1.875e-11})^{1/3} = \qty{64.4}{km}$;
jato $3.628\times64.4 = \qty{233}{km}$.

\textbf{19.} $1.163\times29.4 - A\times\qty{233}{km} = 34.2 - 1.1 =
\qty{33.0}{\sverdrup}$; o excesso, \qty{16}{\percent} do transporte
interior, retorna para o sul na contracorrente logo a leste do jato.

\textbf{20.} $c = \sqrt{9.81\times4000} = \qty{198}{m/s}$; $L_R =
198/\num{8.37e-5} = \qty{2370}{km}$.

\textbf{21.} $198\times12.4206\times3600 = \qty{8860}{km}$.

\textbf{22.} Perímetro $2(6000 + 3336) = \qty{18670}{km}$, \qty{26.2}{h},
2.1 M$_2$ períodos.

\textbf{23.} O círculo tem $2\pi a\cos\ang{35} = \qty{32790}{km}$, percorrido
em um dia lunar, \qty{24.84}{h}: \qty{367}{m/s}, o que exige $c^2/g =
\qty{13.7}{km}$ de água.

\textbf{24.} Com um raio de Rossby comparável à largura da bacia
e o comprimento de onda da M$_2$ comparável ao comprimento da bacia, cerca de um
\omterm{def:b3:ocean-dynamics:amphidrome}{ponto anfidrômico}, perto do centro da bacia; a preamar gira no sentido anti-horário
em torno dele, mantendo a costa à sua direita.

\textbf{25.} $\tfrac34\times0.357\times\cos^2\ang{35} = \qty{0.18}{m}$.

\textbf{26.} A corrente de contorno oeste retorna para o norte os
\textbf{\qty[text-series-to-math=true]{29}{\sverdrup}} que o interior
transporta para o sul (\qty{29.4}{\sverdrup} pelo cálculo).
\end{solution}
